% 量子力学1 量子力学1-17 問題12（134-propagator。当時は 量子力学1-8 問題8） 解いた日 2026-11-08
%% kotae
\toi{1}{◎}{}
\begin{align*}
K&=\int_{-\infty}^{\infty}\bra{x}\exp\Bigl(-\frac{\ii t}{\hbar}\hat{H}\Bigr)\ket{k}\braket{k}{x'}\dd k\\
&=\int_{-\infty}^{\infty}\bra{x}\sum_{n=0}^{\infty}\frac{1}{n!}\Bigl(\frac{-\ii t}{2m\hbar}\Bigr)^{n}\hat{p}^{2n}\ket{k}\\
&\qquad\times\frac{1}{\sqrt{2\pi}}\ee^{-\ii kx'}\dd k\\
&=\int_{-\infty}^{\infty}\bra{x}\sum_{n=0}^{\infty}\frac{1}{n!}\Bigl(\frac{-\ii t}{2m\hbar}\Bigr)^{n}(\hbar k)^{2n}\ket{k}\\
&\qquad\times\frac{1}{\sqrt{2\pi}}\ee^{-\ii kx'}\dd k\\
&=\int_{-\infty}^{\infty}\exp\Bigl(\frac{-\ii\hbar tk^{2}}{2m}\Bigr)\cdot\frac{1}{2\pi}\ee^{\ii k(x-x')}\dd k\\
&=\frac{1}{2\pi}\int_{-\infty}^{\infty}\exp\Bigl(\frac{-\ii\hbar tk^{2}}{2m}+\ii k(x-x')\Bigr)\dd k
\end{align*}
$a=\ii\hbar t/(2m)$，$b=\ii(x-x')$
\begin{align*}
&=\frac{1}{2\pi}\sqrt{\frac{2\pi m}{\ii\hbar t}}\exp\biggl[\frac{\ii^{2}(x-x')^{2}}{4\cdot\frac{\ii\hbar t}{2m}}\biggr]\\
&=\sqrt{\frac{m}{2\pi\ii\hbar t}}\exp\Bigl[\frac{\ii m(x-x')^{2}}{2\hbar t}\Bigr]
\end{align*}
\toi{2}{◎}{}
\begin{align*}
\psi(x,t)&=\int\bra{x}\ee^{-\ii\hat{H}t/\hbar}\ket{x'}\braket{x'}{\psi(0)}\dd x'\\
&=\bra{x}\ee^{-\ii\hat{H}t/\hbar}\ket{\psi(0)}\\
&=\braket{x}{\psi(t)}
\end{align*}
より $\psi(x,t)$ が得られる。
\[ \ii\hbar\pdv{}{t}K(x,t;x')=\hat{H}K(x,t;x') \]
\[ C=\sqrt{\frac{m}{2\pi\ii\hbar}},\qquad d=\frac{\ii m(x-x')^{2}}{2\hbar} \]
\begin{align*}
\pdv{K}{t}&=\pdv{}{t}\bigl(Ct^{-1/2}\exp(dt^{-1})\bigr)\\
&=-\frac{1}{2}Ct^{-3/2}\exp(dt^{-1})\\
&\qquad+Ct^{-1/2}\cdot(-dt^{-2})\exp(dt^{-1})\\
&=\Bigl(-\frac{1}{2}t-d\Bigr)Ct^{-5/2}\exp(dt^{-1})
\end{align*}
\[ \hat{H}K=\frac{\hat{p}^{2}}{2m}K=-\frac{\hbar^{2}}{2m}\frac{\dd^{2}}{\dd x^{2}}K \]
\begin{align*}
\frac{\dd}{\dd x}K&=\sqrt{\frac{m}{2\pi\ii\hbar t}}\cdot\frac{\ii m\cdot2(x-x')}{2\hbar t}\\
&\qquad\times\exp\Bigl[\frac{\ii m(x-x')^{2}}{2\hbar t}\Bigr]
\end{align*}
\begin{align*}
\frac{\dd^{2}}{\dd x^{2}}K&=Ct^{-1/2}\biggl(\frac{\ii m}{\hbar t}\exp\Bigl[\frac{\ii m(x-x')^{2}}{2\hbar t}\Bigr]\\
&\qquad+\Bigl(\frac{\ii m(x-x')}{\hbar t}\Bigr)^{2}\exp\Bigl[\frac{\ii m(x-x')^{2}}{2\hbar t}\Bigr]\biggr)\\
&=Ct^{-1/2}\Bigl(\frac{\ii m}{\hbar t}-\frac{m^{2}(x-x')^{2}}{\hbar^{2}t^{2}}\Bigr)\exp(dt^{-1})
\end{align*}
\begin{align*}
\hat{H}K&=Ct^{-5/2}\Bigl(\frac{-\hbar^{2}t^{2}\ii m}{2m\hbar t}-\frac{-\hbar^{2}t^{2}m^{2}(x-x')^{2}}{2m\hbar^{2}t^{2}}\Bigr)\\
&\qquad\times\exp(dt^{-1})\\
&=Ct^{-5/2}\Bigl(\frac{-\ii\hbar t}{2}+\frac{m(x-x')^{2}}{2}\Bigr)\exp(dt^{-1})\quad\cdots\text{①}
\end{align*}
\begin{align*}
\ii\hbar\pdv{K}{t}&=Ct^{-5/2}\Bigl(\frac{-\ii\hbar t}{2}-\frac{\ii m(x-x')^{2}}{2\hbar}\cdot\ii\hbar\Bigr)\\
&\qquad\times\exp(dt^{-1})\\
&=Ct^{-5/2}\Bigl(\frac{-\ii\hbar t}{2}+\frac{m(x-x')^{2}}{2}\Bigr)\exp(dt^{-1})\quad\cdots\text{②}
\end{align*}
① $=$ ② より $\ii\hbar\,\partial K/\partial t=\hat{H}K$
\toi{3}{◎}{}
\begin{align*}
&\int K\ee^{\ii k_{0}x'}\dd x'\\
&=\sqrt{\frac{m}{2\pi\ii\hbar t}}\int_{-\infty}^{\infty}\exp\Bigl[\frac{\ii m}{2\hbar t}(x'-x)^{2}+\ii k_{0}x'\Bigr]\dd x'
\end{align*}
$x'-x=z$
\begin{align*}
&=\sqrt{\frac{m}{2\pi\ii\hbar t}}\int_{-\infty}^{\infty}\exp\Bigl[\frac{\ii m}{2\hbar t}z^{2}+\ii k_{0}(x+z)\Bigr]\dd z
\end{align*}
$a=-\ii m/(2\hbar t)$，$b=\ii k_{0}$
\begin{align*}
&=\sqrt{\frac{m}{2\pi\ii\hbar t}}\ee^{\ii k_{0}x}\int_{-\infty}^{\infty}\exp\Bigl[\frac{\ii m}{2\hbar t}z^{2}+\ii k_{0}z\Bigr]\dd z\\
&=\sqrt{\frac{m}{2\pi\ii\hbar t}}\ee^{\ii k_{0}x}\cdot\sqrt{-\frac{2\pi\hbar t}{\ii m}}\\
&\qquad\times\exp\biggl[\frac{(\ii k_{0})^{2}}{4\cdot\frac{-\ii m}{2\hbar t}}\biggr]\\
&=\exp\Bigl[\ii k_{0}x-\frac{\ii\hbar k_{0}^{2}t}{2m}\Bigr]\\
&=\exp[\ii k_{0}x]\exp\Bigl[-\ii\Bigl(\frac{\hbar^{2}k_{0}^{2}}{2m}\Bigr)t/\hbar\Bigr]
\end{align*}
エネルギー $\hbar^{2}k_{0}^{2}/(2m)$ の定常状態
\toi{4}{▲}{方針}
\begin{align*}
\abs{K}^{2}&=\sqrt{\frac{-m}{2\pi\ii\hbar t}}\cdot\sqrt{\frac{m}{2\pi\ii\hbar t}}\\
&\qquad\times\exp\Bigl[\frac{-\ii m(x-x')^{2}}{2\hbar t}+\frac{\ii m(x-x')^{2}}{2\hbar t}\Bigr]\\
&=\frac{m}{2\pi\hbar t}
\end{align*}
$x$ によらない。

\ayamari{エネルギーが一定 $\to$ 運動量が一意に定まる。}{運動量の分布は一様}不確定性原理 $\Delta x\,\Delta p\ge\hbar/2$ より，\ayamari{運動量が決まれば}{運動量が定まらないので}位置の確率は全空間に広がる。
\naze{時刻 $0$ の状態 $\ket{x'}$ は位置が確定した状態で，エネルギーの固有状態でも運動量の固有状態でもない。その運動量表示は $\braket{k}{x'}=\ee^{-\ii kx'}/\sqrt{2\pi}$ で，運動量の分布 $\abs{\braket{k}{x'}}^{2}=1/(2\pi)$ はすべての $k$ で同じである。$\Delta x\to0$ なので $\Delta p\to\infty$ で，不確定性関係の向きが逆である。また運動量が定まった平面波は時刻 $0$ から全空間に広がっているので，局在した粒子が広がる説明にもならない。正しくは，どれほど大きな運動量の成分も同じ重みで含まれ，大きな運動量の成分ほど速く（速さ $\hbar k/m$）遠くへ進むので，どれほど短い時間の後にも全空間に広がる。}
%% kizuki
\item 位置の固有ケット $\ket{x'}$ の運動量表示 $\braket{k}{x'}$ は絶対値が $k$ によらない。位置の不確定さが $0$ の極限では運動量の不確定さは無限大で，運動量の分布は一様になる（デルタ関数のフーリエ変換が定数になることと同じ）。自由粒子では運動量 $\hbar k$ の成分は速さ $\hbar k/m$ で進み，$\abs{k}$ に上限が無いので，どれほど短い $t>0$ でもどれほど遠くまで届く成分がある。
\item \mae{第7章 問題3 (4)・11-04，第7章 問題4 (1)・11-07} 運動量の分布を $\abs{\braket{k}{\psi}}^{2}$ で考えていない（前回は井戸のエネルギーの固有状態について運動量で何がいえるかを答えず，$\abs{\tilde{\psi}(k)}^{2}\dd k$ が確率と読める理由を書かなかった）。「エネルギーが定まる」と「運動量が定まる」は別のことである。説明の問題では，まず問題の状態が何の固有状態か（または何の固有状態でもないか）を確かめ，「運動量の分布」と問われたら $\braket{k}{\psi}$ を書いてその絶対値の2乗が $k$ でどう変わるかを見る。
